\documentclass[12pt]{article}
\usepackage[margin=0.6in,top=0.85in,bottom=0.55in]{geometry}
\usepackage{lmodern}
\usepackage{microtype}
\usepackage{fancyhdr}
\usepackage{titlesec}
\usepackage{enumitem}
\usepackage{array}
\usepackage{lastpage}
\usepackage[hidelinks]{hyperref}
\setlength{\parindent}{0pt}
\setlength{\parskip}{0.2em}
\setlength{\headheight}{26pt}
\titleformat{\section}{\large\bfseries\scshape}{}{0pt}{}
\pagestyle{fancy}
\fancyhf{}
\fancyhead[C]{\footnotesize Student Name: \hspace{1.65in} \quad Assignment ID: U1F3LD \quad Page \thepage\ of \pageref{LastPage}}
\renewcommand{\headrulewidth}{0.5pt}
\renewcommand{\footrulewidth}{0pt}
\newcommand{\answerbox}[2]{\par\vspace{0.08em}\noindent\textbf{#1}\par\vspace{0.03em}\noindent\fbox{\begin{minipage}[t][#2]{\dimexpr\textwidth-2\fboxsep-2\fboxrule}\rule{\textwidth}{0pt}\vspace{#2}\end{minipage}}\par\vspace{0.16em}}
\newcommand{\cluebox}[1]{\par\vspace{0.12em}\noindent\begin{minipage}{\textwidth}\itshape #1\end{minipage}\par\vspace{0.12em}}
\begin{document}
\begin{center}
{\LARGE\bfseries Logic Dungeon III}\par\vspace{0.04em}
{\large\scshape The Court of the False Oracle}\par\vspace{0.08em}
\rule{0.78\textwidth}{0.5pt}\par
\end{center}
\textbf{Purpose:} Use the logic tools from Weeks 7--9 to decide which authorization, if any, is best supported as the next step for twenty adventurers.

\section*{The Court}
The company has reached the Upper Archive by Rootbridge. In the \textbf{Court of the Oracle}, three institutional offices issue final-sounding rulings when records are incomplete. Nothing here is supernatural. The danger is procedure that looks like proof.

\begin{itemize}[leftmargin=1.5em,itemsep=0.06em]
\item \textbf{Charge --- Advocate of the Bench:} duty and loyalty rulings.
\item \textbf{Passage --- Master of Ways:} route and advance rulings.
\item \textbf{Seal --- Keeper of Marks:} which documents ``count'' as binding.
\end{itemize}

A petition is open: how may a full company of twenty proceed deeper into the Archive? Three Reply Tablets already sit on the bench. A ceremonial door waits for finished Court authorization. Behind a dusty hanging is a grated \textbf{Records Service Corridor} used by clerks---no seals, no Echo Tablet, easy to dismiss as irregular.

\cluebox{\textbf{Rule:} Diagnose each argument. Record what it establishes, what it does not, and what limited action it justifies. Prefer \emph{independent support} over finished form. The goal is the \emph{best-supported next step}, not a guarantee and not obedience to the loudest office.}

\textbf{Tools:} fallacy as appearance of proof; false claim vs.\ bad inference; emotional/rhetorical force vs.\ logical force; fallacy autopsy; ambiguous middle / full-argument equivocation; fix-and-retest; circular reasoning / begging the question; loaded definition; independent support; repair. Spiral: follow-test, form vs.\ fact, stable terms.
\newpage
\section*{Part 1: Charge --- Danger Station}
The Advocate of the Bench strikes the Echo frame and announces:

\cluebox{``Anyone who delays for more proof is against the Archive's work. You asked for more proof before moving all twenty. Therefore you are against the Archive's work. A true courier accepts a finished Court ruling. Accept the ruling now.''}

\textbf{Question 1.1}\\[-0.15em]
Perform a fallacy autopsy: (1) conclusion, (2) what gives the argument its appearance of proof, (3) the main defect, and (4) a repaired or narrowed claim.
\answerbox{ANSWER LABEL: Part 1 Q1}{2.15in}

\textbf{Question 1.2}\\[-0.15em]
What may the company record after Q1? Does Charge establish a route, a safety fact, both, or neither? What limited action is justified?
\answerbox{ANSWER LABEL: Part 1 Q2}{1.55in}

\newpage
\section*{Part 2: Passage --- Danger Station}
The Master of Ways points to the ceremonial door:

\cluebox{``All sound advances are advances the company should take. The threshold under the ceremonial door is sound---masons found no crack in the first stone. Therefore the ceremonial advance is a sound advance the company should take.''}

\textbf{Question 2.1}\\[-0.15em]
Identify the shifting term. State both meanings, rewrite the argument with one stable meaning, and decide whether the original conclusion still follows.
\answerbox{ANSWER LABEL: Part 2 Q1}{2.2in}

\textbf{Question 2.2}\\[-0.15em]
Separate any useful fact inside the bad argument from what remains unproved about full-company advance. What limited action is justified at the ceremonial threshold?
\answerbox{ANSWER LABEL: Part 2 Q2}{1.55in}

\newpage
\section*{Part 3: Seal --- Danger Station}
The Keeper of Marks holds up a waxed packet labeled \emph{Advance Warrant}:

\cluebox{``This packet is binding because it is an authentic Archive issue. It is an authentic Archive issue because it bears the current binding mark. The current binding mark is valid because binding Archive issues carry it. Therefore the company must treat this packet's order as settled proof of the correct advance.''}

\textbf{Question 3.1}\\[-0.15em]
Explain where the argument is circular (what is assumed that still needed proof). Identify any loaded definition. State what independent support would be needed.
\answerbox{ANSWER LABEL: Part 3 Q1}{2.15in}

\textbf{Question 3.2}\\[-0.15em]
What may the company preserve about the packet, and what diagnosis should they give Seal's authorization right now: accept for all twenty, reject as impossible, or provisional?
\answerbox{ANSWER LABEL: Part 3 Q2}{1.55in}

\newpage
\section*{Part 4: Two Failures on One Bench}
Two short notes are pinned beside the Echo frame.

\cluebox{\textbf{Note A:} ``All side corridors are collapse traps. The Records Service Corridor is a side corridor. Therefore it is a collapse trap.'' A Warden scrap in the same tray says only: ``Service corridor last inspected for clearance; no collapse recorded in the log years we still hold.''\\[0.25em]
\textbf{Note B:} ``Two earlier parties left through a service grate. Therefore the Service Corridor is the authorized Court advance for all companies.''}

\textbf{Question 4.1}\\[-0.15em]
For Note A, is the main failure a false claim, a bad inference, or both? Explain. For Note B, name the main failure and what it does not prove.
\answerbox{ANSWER LABEL: Part 4 Q1}{2.0in}

\textbf{Question 4.2}\\[-0.15em]
After Part 4 alone, may the company commit all twenty to the Service Corridor? What limited record should they keep about it?
\answerbox{ANSWER LABEL: Part 4 Q2}{1.55in}

\newpage
\section*{Part 5: Records That Change the Field}
While no one has yet filled the ceremonial seal-cups, the company opens three short clerical records from the seal cabinet and the Warden's desk.

\cluebox{\textbf{Record R --- Seal Cabinet Card:} ``The die called `current binding mark' was cut after the original ring was lost. Several packets were restamped for tray order. A restamp proves the packet was refiled, not that its order is independently true or safe for a full company.''\\[0.2em]
\textbf{Record S --- Passage Stock List:} ``Many Passage Reply Tablets reuse the word \emph{sound} in standing formulas. Formula series P-12 uses \emph{sound} for `no crack observed on first stone' and, in a later sentence, for `approved full advance.'''\\[0.2em]
\textbf{Record T --- Echo Standing Order:} ``An Echo restatement may repeat a ruling for the room. Echo confirmation is not additional evidence.''}

\textbf{Question 5.1}\\[-0.15em]
How does Record R change Seal's fact status without needing magic or new prophecy? Give Seal's present diagnosis and route consequence.
\answerbox{ANSWER LABEL: Part 5 Q1}{1.75in}

\textbf{Question 5.2}\\[-0.15em]
How do Records S and T strengthen the Part 2 and Part 1 diagnoses? What should the company refuse to treat as independent support?
\answerbox{ANSWER LABEL: Part 5 Q2}{1.75in}

\newpage
\section*{Part 6: The Service Ledger}
Behind the dusty hanging, a retreat cord snags a panel. The company recovers the Warden's corridor book and a maintenance comparison.

\cluebox{\textbf{Record U --- Service Inspection:} ``Records Service Corridor: iron grate sound; floor dry; clearance height for packs; inspected for full clerk crews (up to two dozen) moving sealed trays. Hazards limited to low steps and one narrow turn. No ceremonial authorization required for service traffic.''\\[0.2em]
\textbf{Record V --- Stack Access Note:} ``This corridor is the clerks' working route from Court to the next archive galleries. It does not open every later vault. It does provide a controlled, currently inspected path for a full working party.''\\[0.2em]
\textbf{Record W --- Charge Margin:} ``Loyalty language on Charge tablets is stock formula C-4. It does not report a route inspection.''}

A non-circular case is now available:

\cluebox{Routes currently inspected for full working crews, with dry floor and pack clearance as confirmed by the Service Inspection, are routes the company may take as a controlled next step into the next archive galleries.\\
The Records Service Corridor is such a route under Records U and V.\\
Therefore the company may take the Records Service Corridor as a controlled next step into the next archive galleries.}

\textbf{Question 6.1}\\[-0.15em]
State the conclusion, the independent support (not a restatement), and why this case is not circular. What does Record W add?
\answerbox{ANSWER LABEL: Part 6 Q1}{1.85in}

\textbf{Question 6.2}\\[-0.15em]
What new decision is justified now that was not justified in Part 4? Distinguish ``best-supported under the records'' from ``guaranteed free of every hazard.''
\answerbox{ANSWER LABEL: Part 6 Q2}{1.35in}

\newpage
\section*{Part 7: Court Evidence Chart}
\renewcommand{\arraystretch}{1.22}
\begin{tabular}{|p{0.16\textwidth}|p{0.2\textwidth}|p{0.2\textwidth}|p{0.18\textwidth}|p{0.16\textwidth}|}
\hline
\textbf{Option} & \textbf{Appearance of proof} & \textbf{Main defect} & \textbf{Independent support?} & \textbf{Diagnosis}\\\\ \hline
Charge & & & & \\\\[0.34in] \hline
Passage & & & & \\\\[0.34in] \hline
Seal & & & & \\\\[0.34in] \hline
Service Corridor & & & & \\\\[0.34in] \hline
\end{tabular}

\textbf{Question 7.1}\\[-0.15em]
Complete the chart. Then identify: (a) the argument that mainly uses pressure instead of proof, (b) the argument that fails because a middle term shifts, (c) the circular authorization, and (d) the case with independent support.
\answerbox{ANSWER LABEL: Part 7 Q1}{1.2in}

\textbf{Question 7.2}\\[-0.15em]
Which option is best supported as the next step for the twenty? Use at least three pieces of evidence from different parts and explain why each rejected option has weaker support.
\answerbox{ANSWER LABEL: Part 7 Q2}{1.25in}

\newpage
\section*{Part 8: The Company's Decision}
The Advocate declares: ``The corridor looks irregular, so Charge was right to treat hesitation as disloyalty. Seal still has a stamp, so the packet must settle the advance. A finished Court form settles everything.''

\textbf{Question 8.1}\\[-0.15em]
Separate and repair those claims. State what the Service Corridor case establishes under Records U and V, what remains unguaranteed, and why appearance, stamps, and finished form are not enough by themselves.
\answerbox{ANSWER LABEL: Part 8 Q1}{1.95in}

\textbf{Question 8.2}\\[-0.15em]
Give the company's final movement plan for the best-supported next step. Include one check, one repair or retreat provision, and 2--3 sentences on why independent support mattered more than Court finality here.
\answerbox{ANSWER LABEL: Part 8 Q2}{1.95in}

\end{document}
